\section{Comparison to other BMA software packages}\label{BmaGretlSec4}
Several Bayesian model averaging software packages have been developed by many
researchers%
\footnote{Useful information about Bayesian model averaging software
	packages are available at link
	\url{http://bms.zeugner.eu/resources/}.}.  The most popular in the
literature are \proglang{R} packages.  A recent and comprehensive overview of
Bayesian model averaging software in \proglang{R} is given in
\cite{AminiParmeter2011}. They outline the following BMA packages:
\pkg{BMS} by \cite{Zeugner_BMS}, \pkg{BAS} by
\cite{ClydeLittmanGhosh_BAS} and \pkg{BMA} by
\cite{RafteryHoetingVolinskyPainterYeung_BMA}.  According to their
research the \pkg{BMS} package provides the most accurate results and
it is also the most flexible with numerous options.
Table~\ref{tab:features_BMA_software} presents the main features of
the above mentioned packages.  Furthermore we compare our \pkg{BMA} package
to the \pkg{FLSBMA} program developed in \proglang{Fortran~77} by
\cite{FernandezLeySteel2001b}.

\begin{table}[ht!]
	\caption{Main features of the selected Bayesian model averaging software packages.
	Model sampling: Type of model sampling algorithms.
	Model priors: Number of available model priors.
	Zellner's g-priors: Number of available Zellner's factors of proportionality.
	Jointness analysis: Information if the package provides jointness analysis.
	Jointness measures: Number of available jointness measures.
	Out-of-sample forecast: Information about which package can be used for out-of-sample forecasting.
	No.~of burn-in draws.
	No.~of draws and no.~of top ranked models: Number of user-defined parameters in the sampling algorithm.
	Graphical user interface: Presence of a GUI in the package.}
	\centering
	\begin{tabular}{@{\extracolsep{-1pt}}lccccc}\hline
		Main features			&\proglang{R}-\pkg{BMS}			& \proglang{R}-\pkg{BMA}		& \proglang{R}-\pkg{BAS}	    &\pkg{FLSBMA}	      &\pkg{gretl}-\pkg{BMA}\\\hline
		\\[-10pt]
		\multirow{2}{*}{Model sampling}&\multirow{2}{*}{$\mbox{MC}^3$}		&\multirow{2}{*}{BIC}	&BAS, AMCMC,	    &\multirow{2}{*}{$\mbox{MC}^3$}     &\multirow{2}{*}{$\mbox{MC}^3$}\\
		&				&			&BAS+MCMC	    &			      &\\  
		Model priors			&5				&1			&3		    &2			      &3\\
		Zellner's g-priors	 	&6				&1			&9		    &9			      &5\\
		Jointness analysis	 	&				&			&		    &\checkmark		      &\checkmark\\
		Jointness measures     	&				&			&		    &1			      &7\\
		Out-of-sample forecast 	&\checkmark			&			&\checkmark         &\checkmark		      &\checkmark\\
		No.\ of burn-in draws		&\checkmark			&			&\checkmark         &\checkmark		      &\checkmark\\
		No.\ of draws	 		&\checkmark			&\checkmark		&\checkmark         &\checkmark		      &\checkmark\\
		No.\ of top ranked models  	&\checkmark			&			&	            &\checkmark		      &\checkmark\\
		Graphical user interface	&				&			&		    &			      &\checkmark\\\hline
	\end{tabular}
	\label{tab:features_BMA_software}
\end{table}

The most flexible package according to the model sampling category is
\proglang{R}-\pkg{BAS} which contains three sampling algorithms: Bayesian
adaptive sampling (BAS), adaptive MCMC (AMCMC) and Bayesian adaptive
sampling with MCMC for the initial replications (BAS+MCMC).  All
packages offer different model priors. The most limited is
\proglang{R}-\pkg{BMA} with only a fixed model prior, while the most
comprehensive is \proglang{R}-\pkg{BMS} which contains the following model
priors: fixed, uniform, binomial, binomial-beta and custom
(user-controlled prior probabilities).  Furthermore, the number of
g-priors may differ depending on the package starting from one for
\proglang{R}-\pkg{BMA}, to nine for \proglang{R}-\pkg{BAS} and \pkg{FLSBMA}.

\pkg{gretl}-\pkg{BMA} and \pkg{FLSBMA} are the only packages which offer
Bayesian model averaging together with accompanying jointness
statistics.  \pkg{FLSBMA} contains one jointness measure (Ley-Steel)
while \pkg{gretl}-\pkg{BMA} contains seven measures.
The estimation of out-of-sample forecasts is
possible in almost all packages excluding \proglang{R}-\pkg{BMA}, where this
option is not available.  Three out of five packages enable full
control of the replication algorithm.  In \proglang{R}-\pkg{BMS}, \pkg{FLSBMA}
and \pkg{gretl}-\pkg{BMA}, the user can set the number of burn-in draws,
replication draws and detailed information about top ranked models.
Unlike other packages \pkg{gretl}-\pkg{BMA} offers an easy-to-use GUI for
the Bayesian model averaging approach.

In order to compare run times of the above mentioned packages we used
the FLS dataset with the following options: benchmark g-prior, 1000000
burn-in draws, 2000000 iterations.  Model prior was set to uniform. We
considered two scenarios for our experiment: with or without jointness
analysis.

\begin{table}[H]
	\caption{Run time in seconds of Bayesian model averaging software packages for FLS data.}
	\centering
	\begin{tabular}{lccccc}\hline
		Run time &\proglang{R}-\pkg{BMS} &\proglang{R}-\pkg{BMA} &\proglang{R}-\pkg{BAS} &\pkg{FLSBMA} &\pkg{gretl}-\pkg{BMA}\\\hline
		\\[-10pt]
		Without jointness analysis	&277&87 &53 &	&66\\
		With jointness analysis		&	&	&	&612&73\\\hline
	\end{tabular}	
	\label{tab:FLS_packages_timings}
\end{table}
